% --- Archon physics formalization source begin ---
% archon:physics
% archon:covers PhyXMiniProblems/problem_phyx_mini_0047.lean
% archon:source-report reports/phyx_mini/problem_phyx_mini_0047.source.json

\chapter{Physics problem phyx\_mini\_0047}
\label{ch:PhyXMiniProblems_problem_phyx_mini_0047}

\paragraph{Problem source.}
\#\# Physical scenario

You pass 633\textbackslash{} \textbackslash{}mathrm\{nm\} laser light through a narrow slit and observe the diffraction pattern on a screen 6.0\textbackslash{} \textbackslash{}mathrm\{m\} away. The distance on the screen between the centers of the first minima on either side of the central bright fringe is 32\textbackslash{} \textbackslash{}mathrm\{mm\}.

\#\# Question

How wide is the slit?

\#\# Answer choices

- A: A: 643nm

- B: B: 633nm

- C: C: 639nm

- D: D: 533nm

\#\# Figure caption (auxiliary; use the image as primary evidence)

The image illustrates a diffraction setup. On the left, there is a barrier with a single slit, labeled "Slit width = ?" indicating that the slit width is unknown. A light wave passes through this slit, creating a diffraction pattern on the screen to the right.

The screen is positioned at a distance of 6.0 meters from the slit, as indicated by the label "x = 6.0 m." On the screen, a series of light and dark bands, known as interference fringes, are visible. The central band is the brightest.

The spacing between the fringes on the screen is marked as "32 mm," measured vertically. The axes x and y are labeled at the top right of the screen, with x being horizontal and y vertical.

Overall, the diagram depicts an experimental setup to determine the width of the slit using the interference pattern produced on the screen.

\#\# Dataset metadata

Wave Optics; Physical Model Grounding Reasoning, Multi-Formula Reasoning

\paragraph{Recorded answer/context.}
B: B: 633nm

\paragraph{Figure/image path.}
/root/proposal\_for\_physic/science-mango/phyx\_mini\_run/phyx\_data/test\_image/47.png

\paragraph{Formalization target.}
create a compiling Lean file with sorry bodies at `PhyXMiniProblems/problem\_phyx\_mini\_0047.lean`. The Lean declarations must preserve the physical quantities, dimensions or dimensional roles, figure labels, governing-law hypotheses, and final relation expressed by this problem.
Use Mathlib/Physlib names found through LeanExplore where available. If a physics API is missing, introduce faithful local abstractions rather than scalar placeholder aliases.

\begin{theorem}[Physics formalization target]
\label{thm:physics:phyx_mini_0047:target}
\lean{PhyXMiniProblems.ProblemPhyXMini0047.problem_phyx_mini_0047}
\uses{def:physics:phyx-mini-0047:phyxminiproblems-problemphyxmini0047-lengthinnanometers, def:physics:phyx-mini-0047:phyxminiproblems-problemphyxmini0047-singleslitapparatus, def:physics:phyx-mini-0047:phyxminiproblems-problemphyxmini0047-firstminimumgeometry, def:physics:phyx-mini-0047:phyxminiproblems-problemphyxmini0047-matchesproblemandfigurereadouts, def:physics:phyx-mini-0047:phyxminiproblems-problemphyxmini0047-hasphysicalparameters, def:physics:phyx-mini-0047:phyxminiproblems-problemphyxmini0047-usesfraunhofersingleslitmodel, def:physics:phyx-mini-0047:phyxminiproblems-problemphyxmini0047-satisfiesfirstminimumdiffractionlaws}
The assigned autoformalize agent should translate this physics problem into Lean declarations in the covered file, with theorem and lemma proof bodies written as `by sorry`.
\end{theorem}
\begin{proof}
This is an autoformalization task, not a proof task. Produce statements that can later be proved by the physics prover without weakening the physical model.
\end{proof}

% --- Archon declaration topology begin ---
\section{Declaration topology}
\begin{definition}[Optical Length]
  \label{def:physics:phyx-mini-0047:phyxminiproblems-problemphyxmini0047-opticallength}
  \lean{PhyXMiniProblems.ProblemPhyXMini0047.OpticalLength}
  A signed, unit-independent physical quantity carrying length dimension
\end{definition}
\begin{proof}
  This is the declared model definition of Optical Length.
\end{proof}
\begin{definition}[length Readout]
  \label{def:physics:phyx-mini-0047:phyxminiproblems-problemphyxmini0047-lengthreadout}
  \lean{PhyXMiniProblems.ProblemPhyXMini0047.lengthReadout}
  \uses{def:physics:phyx-mini-0047:phyxminiproblems-problemphyxmini0047-opticallength}
  Read a physical length as a real scalar in the selected length unit
\end{definition}
\begin{proof}
  This is the declared model definition of length Readout.
\end{proof}
\begin{definition}[length In Meters]
  \label{def:physics:phyx-mini-0047:phyxminiproblems-problemphyxmini0047-lengthinmeters}
  \lean{PhyXMiniProblems.ProblemPhyXMini0047.lengthInMeters}
  \uses{def:physics:phyx-mini-0047:phyxminiproblems-problemphyxmini0047-opticallength, def:physics:phyx-mini-0047:phyxminiproblems-problemphyxmini0047-lengthreadout}
  Metre readout for axial geometry and positivity conditions
\end{definition}
\begin{proof}
  This is the declared model definition of length In Meters.
\end{proof}
\begin{definition}[length In Millimeters]
  \label{def:physics:phyx-mini-0047:phyxminiproblems-problemphyxmini0047-lengthinmillimeters}
  \lean{PhyXMiniProblems.ProblemPhyXMini0047.lengthInMillimeters}
  \uses{def:physics:phyx-mini-0047:phyxminiproblems-problemphyxmini0047-opticallength, def:physics:phyx-mini-0047:phyxminiproblems-problemphyxmini0047-lengthreadout}
  Millimetre readout for the screen separation in the source figure
\end{definition}
\begin{proof}
  This is the declared model definition of length In Millimeters.
\end{proof}
\begin{definition}[length In Nanometers]
  \label{def:physics:phyx-mini-0047:phyxminiproblems-problemphyxmini0047-lengthinnanometers}
  \lean{PhyXMiniProblems.ProblemPhyXMini0047.lengthInNanometers}
  \uses{def:physics:phyx-mini-0047:phyxminiproblems-problemphyxmini0047-opticallength, def:physics:phyx-mini-0047:phyxminiproblems-problemphyxmini0047-lengthreadout}
  Nanometre readout for the wavelength, answers, and requested width
\end{definition}
\begin{proof}
  This is the declared model definition of length In Nanometers.
\end{proof}
\begin{definition}[Aperture Geometry]
  \label{def:physics:phyx-mini-0047:phyxminiproblems-problemphyxmini0047-aperturegeometry}
  \lean{PhyXMiniProblems.ProblemPhyXMini0047.ApertureGeometry}
  The aperture geometry selected by the single opening in the barrier
\end{definition}
\begin{proof}
  This is the declared model definition of Aperture Geometry.
\end{proof}
\begin{definition}[Diffraction Regime]
  \label{def:physics:phyx-mini-0047:phyxminiproblems-problemphyxmini0047-diffractionregime}
  \lean{PhyXMiniProblems.ProblemPhyXMini0047.DiffractionRegime}
  The far-field diffraction model used for the first-minimum law
\end{definition}
\begin{proof}
  This is the declared model definition of Diffraction Regime.
\end{proof}
\begin{definition}[Single Slit Apparatus]
  \label{def:physics:phyx-mini-0047:phyxminiproblems-problemphyxmini0047-singleslitapparatus}
  \lean{PhyXMiniProblems.ProblemPhyXMini0047.SingleSlitApparatus}
  \uses{def:physics:phyx-mini-0047:phyxminiproblems-problemphyxmini0047-opticallength, def:physics:phyx-mini-0047:phyxminiproblems-problemphyxmini0047-aperturegeometry, def:physics:phyx-mini-0047:phyxminiproblems-problemphyxmini0047-diffractionregime}
  The monochromatic laser, aperture, and observation screen in the experiment. The unknown physical quantity is slitWidth; no field assigns it a numerical value or defines it from an answer choice
\end{definition}
\begin{proof}
  This is the declared model definition of Single Slit Apparatus.
\end{proof}
\begin{definition}[First Minimum Geometry]
  \label{def:physics:phyx-mini-0047:phyxminiproblems-problemphyxmini0047-firstminimumgeometry}
  \lean{PhyXMiniProblems.ProblemPhyXMini0047.FirstMinimumGeometry}
  \uses{def:physics:phyx-mini-0047:phyxminiproblems-problemphyxmini0047-opticallength}
  Signed screen coordinates and propagation angles for the two first minima. The central-axis coordinate makes the origin drawn at the intersection of the figure's x and y axes explicit
\end{definition}
\begin{proof}
  This is the declared model definition of First Minimum Geometry.
\end{proof}
\begin{definition}[Matches Problem And Figure Readouts]
  \label{def:physics:phyx-mini-0047:phyxminiproblems-problemphyxmini0047-matchesproblemandfigurereadouts}
  \lean{PhyXMiniProblems.ProblemPhyXMini0047.MatchesProblemAndFigureReadouts}
  \uses{def:physics:phyx-mini-0047:phyxminiproblems-problemphyxmini0047-lengthreadout, def:physics:phyx-mini-0047:phyxminiproblems-problemphyxmini0047-lengthinmeters, def:physics:phyx-mini-0047:phyxminiproblems-problemphyxmini0047-lengthinmillimeters, def:physics:phyx-mini-0047:phyxminiproblems-problemphyxmini0047-lengthinnanometers, def:physics:phyx-mini-0047:phyxminiproblems-problemphyxmini0047-singleslitapparatus, def:physics:phyx-mini-0047:phyxminiproblems-problemphyxmini0047-firstminimumgeometry}
  The text and primary-image data: one narrow slit, λ = 633 nm, x = 6.0 m, a centered y = 0 optical axis, and 32 mm between the centers of the upper and lower first minima. This predicate contains no value for the unknown slit width
\end{definition}
\begin{proof}
  This is the declared model definition of Matches Problem And Figure Readouts.
\end{proof}
\begin{definition}[Has Physical Parameters]
  \label{def:physics:phyx-mini-0047:phyxminiproblems-problemphyxmini0047-hasphysicalparameters}
  \lean{PhyXMiniProblems.ProblemPhyXMini0047.HasPhysicalParameters}
  \uses{def:physics:phyx-mini-0047:phyxminiproblems-problemphyxmini0047-lengthinmeters, def:physics:phyx-mini-0047:phyxminiproblems-problemphyxmini0047-singleslitapparatus, def:physics:phyx-mini-0047:phyxminiproblems-problemphyxmini0047-firstminimumgeometry}
  Physical branch conditions for the nondegenerate apparatus. The two rays lie on the signed acute branches shown in the figure. These conditions do not determine the numerical slit width
\end{definition}
\begin{proof}
  This is the declared model definition of Has Physical Parameters.
\end{proof}
\begin{definition}[Uses Fraunhofer Single Slit Model]
  \label{def:physics:phyx-mini-0047:phyxminiproblems-problemphyxmini0047-usesfraunhofersingleslitmodel}
  \lean{PhyXMiniProblems.ProblemPhyXMini0047.UsesFraunhoferSingleSlitModel}
  \uses{def:physics:phyx-mini-0047:phyxminiproblems-problemphyxmini0047-singleslitapparatus}
  The one-slit apparatus is interpreted in the Fraunhofer regime
\end{definition}
\begin{proof}
  This is the declared model definition of Uses Fraunhofer Single Slit Model.
\end{proof}
\begin{definition}[Satisfies First Minimum Diffraction Laws]
  \label{def:physics:phyx-mini-0047:phyxminiproblems-problemphyxmini0047-satisfiesfirstminimumdiffractionlaws}
  \lean{PhyXMiniProblems.ProblemPhyXMini0047.SatisfiesFirstMinimumDiffractionLaws}
  \uses{def:physics:phyx-mini-0047:phyxminiproblems-problemphyxmini0047-lengthreadout, def:physics:phyx-mini-0047:phyxminiproblems-problemphyxmini0047-singleslitapparatus, def:physics:phyx-mini-0047:phyxminiproblems-problemphyxmini0047-firstminimumgeometry}
  Exact first-order Fraunhofer single-slit minima and straight slit-to-screen ray geometry. In any coherent length unit the laws are a sin θ\_upper = λ, a sin θ\_lower = -λ, and y - y₀ = L tan θ. The opposite-angle field records the reflection symmetry of the two first minima about the central optical axis. No field contains the requested width or any displayed answer
\end{definition}
\begin{proof}
  This is the declared model definition of Satisfies First Minimum Diffraction Laws.
\end{proof}
\begin{definition}[Answer Choice]
  \label{def:physics:phyx-mini-0047:phyxminiproblems-problemphyxmini0047-answerchoice}
  \lean{PhyXMiniProblems.ProblemPhyXMini0047.AnswerChoice}
  Labels of the four slit-width answers printed with the dataset item
\end{definition}
\begin{proof}
  This is the declared model definition of Answer Choice.
\end{proof}
\begin{definition}[displayed Slit Width In Nanometers]
  \label{def:physics:phyx-mini-0047:phyxminiproblems-problemphyxmini0047-displayedslitwidthinnanometers}
  \lean{PhyXMiniProblems.ProblemPhyXMini0047.displayedSlitWidthInNanometers}
  \uses{def:physics:phyx-mini-0047:phyxminiproblems-problemphyxmini0047-answerchoice}
  Nanometre value printed beside each answer label
\end{definition}
\begin{proof}
  This is the declared model definition of displayed Slit Width In Nanometers.
\end{proof}
\begin{definition}[recorded Dataset Answer]
  \label{def:physics:phyx-mini-0047:phyxminiproblems-problemphyxmini0047-recordeddatasetanswer}
  \lean{PhyXMiniProblems.ProblemPhyXMini0047.recordedDatasetAnswer}
  \uses{def:physics:phyx-mini-0047:phyxminiproblems-problemphyxmini0047-answerchoice}
  The answer label recorded by the source dataset, retained as metadata only
\end{definition}
\begin{proof}
  This is the declared model definition of recorded Dataset Answer.
\end{proof}
% --- Archon declaration topology end ---
% --- Archon physics formalization source end ---
